\documentclass[11pt]{article}

% Packages
\usepackage{amsmath,amsthm,amssymb}
\usepackage{algorithm}
\usepackage{algorithmic}
\usepackage{geometry}
\usepackage{natbib}
\usepackage{graphicx}
\usepackage{booktabs}
\usepackage{mathtools}
\usepackage{enumitem}
\usepackage{parskip}

\geometry{margin=1in}

% Theorem environments
\newtheorem{theorem}{Theorem}[section]
\newtheorem{lemma}[theorem]{Lemma}
\newtheorem{proposition}[theorem]{Proposition}
\newtheorem{corollary}[theorem]{Corollary}
\theoremstyle{definition}
\newtheorem{definition}[theorem]{Definition}
\newtheorem{example}[theorem]{Example}
\theoremstyle{remark}
\newtheorem{remark}[theorem]{Remark}

% Commands
\newcommand{\R}{\mathbb{R}}
\newcommand{\Z}{\mathbb{Z}}
\newcommand{\N}{\mathbb{N}}
\newcommand{\norm}[1]{\|#1\|}
\newcommand{\abs}[1]{|#1|}
\newcommand{\argmin}{\operatorname{argmin}}
\newcommand{\argmax}{\operatorname{argmax}}
\DeclareMathOperator{\conv}{conv}
\DeclareMathOperator{\rank}{rank}
\DeclareMathOperator{\diag}{diag}
\DeclareMathOperator{\Tr}{Tr}
\DeclareMathOperator{\vol}{vol}

\title{Notes}

\author{Brannon King}

\date{\today}

\begin{document}

\maketitle

\section{Math}

\subsection{NU where N is null space basis}
Let $N$ be a basis for the null space of $A$ where $Ax=b$.
Then $NU$ is also a basis for the null space of $A$ for any unimodular $U$.
Hence, we can choose to not apply $U$ in that substitution position,
which can be valuable if $U$ increases the density of $N$.

Suppose we have some other $U' \in \mathbb{Z}^{n \times n}$. Let $\tilde{A}= AU'^{-1}$.
Then
$\tilde{A} (U' N) = A U'^{-1} U' N = A N = 0$, so we can interprert $U'N$ as null space basis for $AU'^{-1}$.

\subsection{Square root of ellipsoid hession gives round ellipsoid}

\begin{align}
\{x : (x - x_0)^T H (x - x_0) \leq 1 \} = \\
\{x : \| H^{1/2} (x - x_0) \|_2 \leq 1 \} = \\
\{x : \| y \|_2 \leq 1, y = H^{1/2} (x - x_0) \}
\end{align}

\subsection{Rounded ellipsoid is maintained through LLL}
Let $B=s_1 H^{1/2} U$ where $B,U=LLL(s_1 H^{1/2})$ and $s_1$ is some scaling factor.
\begin{align}
    \text{Substitute } x = Uy + x_0 \\
    \{y : \| H^{1/2} U y \|_2 \leq 1 \} \\
    \{y : \| B y / s_1 \|_2 \leq 1 \} \\
    B^T B \approx s_2 I & \text{ (since it should be near orthogonal)} \\
    \{y : \frac{s_2}{s_1} \| y \|_2 \leq 1 \}
\end{align}
The same substitution should work for $H^{-1/2}$.

\subsection{Same works for $U^{-1} sub$}
Let $BU^{-1}=s_1 H^{-1/2}$ where $B,U=LLL(s_1 H^{-1/2})$ and $s_1$ is some scaling factor.
\begin{align}
    \text{Substitute } x = U^{-1}y + x_0 \\
    \{y : \| H^{1/2} U^{-1} y \|_2 \leq 1 \} \\
    \{y : y^T U^{-T} H U^{-1} y  \leq 1 \} \\
    \{y : y^T (U H^{-1} U^T)^{-1} y  \leq 1 \} \\
    \{y : y^T (U H^{-1} U^T)^{-1} y  \leq 1 \} \\
    \{y : y^T (BB^T / s_1^2)^{-1} y  \leq 1 \}
\end{align}
Again, since $B$ is mostly orthogonal, it should work. (The inverse of $I$ is orthogonal.)

\subsection{Orthogonality of constraints is improved by rounding ellipsoid}
For some measure of orthogonality $m_0$, show $m_0(AH) < m_0(A)$:
\begin{align}
    H = A^T D A = Q \Lambda Q \text{ (eigen decomposition)} \\
    \sqrt{H} = D^{1/2} Q \\
    \text{show } m_0(sD^{1/2} Q U) < m_0(sD^{1/2} Q) \\
    ... ?
\end{align}

Possible measures of orthogonality:
\begin{itemize}[noitemsep]
    \item Condition number
    \item $\lambda_{min} / \lambda_{max}$
    \item Orthogonality defect: $\left( \prod_i \|b_i\|_2 \right) / |\det(BB^T)|$
\end{itemize}

\subsection{We can do the substitution with homogenous coordinates}
Semicolon is row separator. Substitute $[x; 1] = U[y; 1]$.\\
Run LLL on $s_1 H^{1/2} [I x_0]$ to get $B, U = LLL(s_1 [H^{1/2} H^{1/2}x_0])$.
\begin{align}
    \{y : \| H^{1/2} [I x_0][x; -1] \|_2 \leq 1 \} \\
    ...
\end{align}

\subsection{How far from orthogonal can B be?}
From ChatGPT:
\begin{align}
    K_2(B)=\sqrt(\lambda_{max} / \lambda_{min}) \leq poly(n)(\delta - 1/4)^{(1-n)/2}
\end{align}

\subsection{Rounding ellipsoid transforms $Ax \leq b$ into isotropic system}
Where isotropic means that the rows average out to be orthogonal.
\begin{align}
    \text{Given } Ax + s = b \\
    H = A^T D A \text{ for some D with positive diagonal} \\
    \text{Substitute } x = H^{-1/2} y + x_0 \\
    ...
\end{align}
Still need to condense this: https://chatgpt.com/share/6932fddd-2850-8000-9e24-66fda3ae905e

\subsection{Nearer facets contribute more to ellipsoid}
The contribution: $\sum_i {n_i n_i^T} / s_i^2$ means that facets with small $s_i$ (nearer the center) contribute more to the ellipsoid shape than facets with large $s_i$ (farther from the center).

\subsection{Orthogonal constraints give orthogonal corner edge vectors}

\subsection{GMI cuts benefit from orthogonal rays}

Use effectiveness measure: $(\pi^T x - \beta) / \| \pi \|$

\begin{align}
    r_i = (B^{-1} N)_i = e^T_i B^{-1} N \\
    \|r_i\| \leq \| B^{-1} \| \| N_i \| \\
    \pi_i = frac(r_i) \leq \|r_i\|
\end{align}

So we minimize the denominator by minimizing $\| r_i \|$.

If A has orthogonal columns, then $A_B^{-1}$ should also:
\begin{align}
    A^T A = I + S \text{ (S small)} \\
    A^{-1} A A A^{-1} = A^{-1} (I + S) A^{-1} \\
    I = A^{-1} A^{-T} + A^{-1} S A^{-T} \\
    A^{-1} A^{-T} = I - A^{-1}  S A^{-T} \\
    A^{-1} A^{-T} \approx I \\
    A^{-1} \approx A^{T}
\end{align}

\subsection{$x_0$ can be made integer as part of the rounder?}
Use homogenous coordinates in LLL reduction:
\begin{align}
    B, U = LLL(s [H^{1/2} x_0; 0 1]) \\
    U = [U_1 p; q r] \\
    x=U_1 y + p \\
    -1 = q^T y + r
\end{align}

\pagebreak

\section{Ideas in progress}

\subsection{The essence of the question}

Let $AHL(A, b)$ be the mechanism from Aardal's diophantine paper. It returns an integer nullspace basis $N$ and an offset $x_1$ such that $Ax=b$ iff $x = Ny + x_1, y \in \mathbb{Z}^{n-m}$.
Also, let $N^{-1}$ be the integer, rectangular inverse of $N$, typically gleaned from repeated $AHL(N, e_j)$.

You are given a problem of this form:\\
$P_1 \coloneqq min~\{ c^T x : Ax=b, l \leq x \leq u, x \in \mathbb{Z}^n \}$.

Let $N, x_1 = AHL(A, b)$ and substitute $x = Ny + x_1$ to get:\\
$P_2 \coloneqq min~\{ c^T (Ny + x_1) : l \leq Ny + x_1 \leq u, y \in \mathbb{Z}^{n-m} \}$.

Suppose that, alternately, we came up with some unimodular $U$ (with integer offset $x_2$) that improves the geometry of $A$.

Let's substitute that in with $x = Uw + x_2$ to get:\\
$P_3 \coloneqq min~\{ c^T (Uw + x_2) : AUw=b-Ax_2, l \leq Uw + x_2 \leq u, w \in \mathbb{Z}^n \}$.

Note that $P_3$ is typically more difficult than $P_1$ because it has more constraints (replacing bounds), and U is generally dense with large numbers (being derived from HNF or SNF or other scale-then-round mechanisms).

Let $N', x_1' = AHL(AU, b - Ax_2)$ and substitute $w = N'y + x_1'$ to get:\\
$P_4 \coloneqq min~\{ c^T (U(N'y + x_1') + x_2) : l \leq U(N'y + x_1') + x_2 \leq u, y \in \mathbb{Z}^{n-m} \}$.

Via the principle in section 1.1, $UN'$ is a null space basis for $AU^{-1}$. I'm not sure how that helps, though. In general, $UN'$ will have worse numerics than $N$ because $U$ is dense with large numbers.

Questions:
\begin{enumerate}
    \item Is there any possible unimodular $U$ that improves the geometry of $A$ enough to make up for the increased difficulty of $P_4$ over $P_2$ (or $P_3$ over $P_1$)?
    \item $N^{-1}$ can be used directly to generate cuts for $P_1$. Can $N'^{-1}$ also be used to generate cuts for $P_1$?
    \item Are the cuts generated by $N'^{-1}$ (on either $P_1$ or $P_3$) stronger than those generated by $N^{-1}$?
\end{enumerate}

\subsection{Substitution without integer}

$N^T H N$ gives a reduced transform; we're now in y. We can do this assuming $x_0 = x_1 = min~\{cx: Ax=b, \epsilon \leq x \leq u-\epsilon, x \in R^n\} $.
Actually, let's revisit this: We know that our ellipsoid is not going to produce integer coefficients. Hence, finding $x_0 - x_1$ integer seems like a waste.
Suppose $N = null(A)$ and $N \not \in \mathbb{Z}$ and $x_0 = x_1$:
First substitution: $min~c(Ny + x_0): A(Ny + x_0)=b, 0 \leq Ny + x_0 \leq u, Ny + x_0 \in \mathbb{Z}$ 
-- actually, what is the performance on this? $min c(Ny + x_0): Ny + x_0 = x, 0 \leq x \leq u, x \in \mathbb{Z}$
$T$ derived from $y N^T H N y$: $T = H^{1/2} N$
Substitute $y = Tz$
Second sub: $min~c(NTz + x_0): 0 \leq NTz + x_0 \leq u, NTz + x_0 \in \mathbb{Z}$ (assumes $Ax_0=b$ and $AN=0$)
Actually, $T$ can't do any better than orthogonalizing $N$, right? Maybe it could do better at orthogonalizing the constraints right at $x_0$. Maybe rounding them does better than orthogonalizing them? seems unlikely.

Results: this fails to eliminate the quality constraint because you still require $NTz + x_0=x$. It doesn't make the problem any smaller; it makes it larger by adding y or z vars that are real.

\subsection{Double substitution with integer}

So I come up with a unimodular transform U based on H.5 @ N and LLL with scaling. The problem is $x_0 \neq x_1$ (unless we're lucky or do work to find $x_0$ near $x_r$):
$$(Ny + x_0 - x_1)^T H (Ny + x_0 - x_1)$$ has to get to form: $$(y - x_2)^T H (y - x_2)$$ . let $x_2 = x_1 - x_0$
$$(N (y - W(x_1 - x_0)))^T H N (y - Wx_2)$$
$$(y - Wx_2)^T N^T H N (y - Wx_2)$$

Now we do a second substitution: $y = Uz + Wx_2$  (not sure if we need the $Wx_2$ term)
Actually, if $Wx_2$ is not integer, we're broken on that substitution (or we need a constraint $Uz + Wx_2 \in \mathbb{Z}$). Assume it's integer:
$$min \{ c(N(Uz + Wx_2) + x_0) : 0 \leq N(Uz + Wx_2) + x_0 \leq u, z \in \mathbb{Z}  \}$$

Let's revisit $x_1 == x_0$: $$min ||x - x^*||_2: Ax=b, x \in \mathbb{Z}^n.$$ Becomes: $$min \{|| x_0 + Ny - x^* ||_2 : y \in \mathbb{Z}^{n-m} \}.$$ Do add an additional constraint: $x_0 + Ny \leq u$ to keep our ellipsoid in bounds.
Use Babai's estimator to get a starting value of $y$ (though it may not be necessary).



\end{document}

